Includes classes, namespaces, and files

For use in conjunction with Matlab in order to solve for Phase Space Density, Diffusion, Convection etc.

Diffusion + convection code, V\+E\+R\+B4\+D.

The code solves the Fokker-\/\+Planck equation with convection terms\+: \[ \frac{\partial f}{\partial t} = v_{\Phi} \frac{\partial f}{\partial \Pi} + v_{R} \frac{\partial f}{\partial R} + \frac{1}{G} \frac{\partial}{\partial L} G D_{LL} \frac{\partial f}{\partial L} + \frac{1}{G} \frac{\partial}{\partial V} G (D_{VV} \frac{\partial f}{\partial V} + D_{VK} \frac{\partial f}{\partial K}) + \frac{1}{G} \frac{\partial}{\partial K} G (D_{KV} \frac{\partial f}{\partial V} + D_{KK} \frac{\partial f}{\partial K}) + Losses * f + Sources \] where f is a Phase Space Density (P\+S\+D); Φ is M\+L\+T angle, R is radial distance from Earth center, $ V≡μ∙(K+0.5)^2 $ is a combination of adiabatic invariants μ and $ K≡ {J}{\sqrt{8μ}} $; L is the third adiabatic invariant; D are bounce-\/averaged diffusion coefficients; $ v_Φ $, $ v_R $ are drift velocities; G is the Jacobian of the transformation from (μ,J,L) to (V,K,L)

Books, required to understand the code\+:
\begin{DoxyItemize}
\item Stroustrup C++
\item Schulz and Lanzerotti, 1974
\end{DoxyItemize}

V\+A\+R\+I\+A\+B\+L\+E\+S for understanding code
\begin{DoxyItemize}
\item P\+S\+D -\/ Phase Space Density
\item P -\/ Phi\+: magnetic local time
\item R -\/ Radial distance actual
\item V -\/ Energy
\item K -\/ Pitch andle
\item L -\/ Radial distance distorted
\item Dxx -\/ Diffusion -\/ subscripts denote with respect to what derivative
\end{DoxyItemize}

Convection is solved using \hyperlink{_convection__1_d___u_l_t_i_m_a_t_e___q_u_i_c_k_e_s_t6_8h}{Convection\+\_\+1\+D\+\_\+\+U\+L\+T\+I\+M\+A\+T\+E\+\_\+\+Q\+U\+I\+C\+K\+E\+S\+T6.\+h} which implements Leonard B\+P (1988) Universal Limiter for transient interpolation modeling of the advective transport equations\+: the U\+L\+T\+I\+M\+A\+T\+E conservative difference scheme, N\+A\+S\+A technical Memorandum 100916 I\+C\+O\+M\+P-\/88-\/11 \href{http://www.hadian.ir/teaching/CompHydr/3.pdf}{\tt http\+://www.\+hadian.\+ir/teaching/\+Comp\+Hydr/3.\+pdf}

Diffusion is solved using \hyperlink{_diffusion__1_d_8cpp_ae15e469cb777a21a2dbd751f1c072618}{Diffusion\+\_\+1\+D()} and \hyperlink{_diffusion__2_d_8cpp_ad5785270d5a28bd7a7b8ca31f206c8ee}{Diffusion\+\_\+2\+D()} for the standard cases.

Three other methods are also used for finding diffusion in 2\+D\+: these can be found in \hyperlink{_diffusion___a_d_i1_8h}{Diffusion\+\_\+\+A\+D\+I1.\+h}, \hyperlink{_diffusion___a_d_i2_8h}{Diffusion\+\_\+\+A\+D\+I2.\+h}, and \hyperlink{_diffusion___a_d_i3_8h}{Diffusion\+\_\+\+A\+D\+I3.\+h}

The parameters.\+ini file created in matlab will specify the inversion method. If \char`\"{}\+Lapack\char`\"{} is specified then \hyperlink{_diffusion__2_d_8cpp_ad5785270d5a28bd7a7b8ca31f206c8ee}{Diffusion\+\_\+2\+D()} will be used. Otherwise 1 of the 3 A\+D\+I methods will be used to inversion. All 3 A\+D\+I methods can use multiple threads while Lapack can only use 1. The current chosen method of inversion is Diffusion\+\_\+\+A\+D\+I3() if \char`\"{}\+A\+D\+I\char`\"{} is specified. It can be changed in \hyperlink{_v_e_r_b4_d___solver_8cpp}{V\+E\+R\+B4\+D\+\_\+\+S\+O\+L\+V\+E\+R.\+cpp}

From current testing all of the A\+D\+I methods produce bad results -\/ namely negative P\+S\+D values. All examples should thus be run with Lapack until the A\+D\+I methods are fixed. Unfortunately Lapack is not compatable with multiple threads so running the examples with Lapack will be much slower than A\+D\+I and not scalable. Newer versions of Lapack may allow multithreading however.

The main solver function can be found in \hyperlink{_v_e_r_b4_d___solver_8cpp}{V\+E\+R\+B4\+D\+\_\+\+Solver.\+cpp} It consists of a series of P\+S\+D calculations done for every time step that is specified.

This will be used in order to take the input data from matlab store them into \hyperlink{class_matrix1_d}{Matrix1\+D}, \hyperlink{class_matrix2_d}{Matrix2\+D}, \hyperlink{class_matrix3_d}{Matrix3\+D}, and \hyperlink{class_matrix4_d}{Matrix4\+D} using \hyperlink{_read_initial_data_8h}{Read\+Initial\+Data.\+h}

The matrix classes listed above are used to store the data in 1,2,3 or 4 dimension. One example is a 4\+D matrix containing P,R,V,K

There is also the standardized \hyperlink{class_matrix_n_d}{Matrix\+N\+D} for any size matrices which is used in conjunction with \hyperlink{class_updatable_matrix}{Updatable\+Matrix} which can be updated with new data.

Most of the computation such as numerically approximating derivatives are done with \hyperlink{_matrix_solver_8h}{Matrix\+Solver.\+h} using \hyperlink{class_calculation_matrix}{Calculation\+Matrix}. These calculation matrices are made up of \hyperlink{class_diag_matrix}{Diag\+Matrix} which is a mapping of an int (which diagonal number) to a 1d matrix of values (for that diagonal).

There is also a \hyperlink{class_parameters}{Parameters} class which holds paramters and their value as defined in the file they came from. These parameters get the values specified in the parameters.\+ini text file and set variables to determine inversion method, number of threads, using .mat files, etc.

This code was built with M\+A\+T\+L\+A\+B capabilities including reading and writing .mat files. In order to switch all files to .mat for increased speed and percision change the use\+\_\+matlab variable in the M\+A\+T\+L\+A\+B Conv\+\_\+\+Dif.\+m file for whichever example you are running to true. If the machine that this code is being run on does not have matlab installed change the M\+A\+T\+L\+A\+B\+\_\+\+C\+A\+P\+A\+B\+L\+E variable in \hyperlink{_matrix_8h}{Matrix.\+h} to false 